\chapter{Programmation linéaire : la méthode du simplexe}
\label{chap:linear-programming}

\section{Optimalité et dualité}

\subsection{Optimalité du problème primal}

Considérons le problème de programmation linéaire (PL) sous sa forme standard :
\begin{equation}
    \min_{x \in \mathbb{R}^n} c^T x \quad \text{sujet à} \quad Ax = b, \quad x \geq 0,
\end{equation}
où $A$ est une matrice $m \times n$ de rang $m$, $c \in \mathbb{R}^n$ est le vecteur de coût et $b \in \mathbb{R}^m$.

Une étape clé est d'identifier une \textbf{solution de base réalisable}. On partitionne la matrice $A$ en $A = [B \mid N]$ où $B$ est une matrice $m \times m$ inversible (la base) et $N$ contient les $n-m$ autres colonnes. De même, on partitionne $x = (x_B, x_N)$ et $c = (c_B, c_N)$.
L'équation $Ax=b$ devient $Bx_B + Nx_N = b$, soit $x_B = B^{-1}b - B^{-1}Nx_N$.

En fixant $x_N = 0$, on obtient la solution de base $x_B = B^{-1}b$. Si $x_B \ge 0$, la solution est dite réalisable.

Le coût réduit ("reduced cost") est défini par le vecteur $s_N$:
\begin{equation}
    s_N = c_N - (B^{-1}N)^T c_B = c_N - N^T \lambda,
\end{equation}
avec $\lambda = B^{-T} c_B$ le vecteur des multiplicateurs (ou variables duales).

\begin{theoreme}[Condition d'optimalité]
    Une solution de base réalisable $(x_B, 0)$ est optimale si et seulement si le coût réduit est positif ou nul, c'est-à-dire $s_N \geq 0$ (pour toutes ses composantes).
\end{theoreme}

\subsection{Le problème dual}

À tout problème primal (P) correspond un problème dual (D). Pour la forme standard ci-dessus, le dual est :
\begin{equation}
    \max_{\lambda \in \mathbb{R}^m} b^T \lambda \quad \text{sujet à} \quad A^T \lambda \leq c.
\end{equation}

Les relations fondamentales entre primal et dual sont :
\begin{itemize}
    \item \textbf{Dualité faible} : Pour tout $x$ réalisable pour (P) et tout $\lambda$ réalisable pour (D), on a $b^T \lambda \leq c^T x$.
    \item \textbf{Dualité forte} : Si (P) admet une solution optimale $x^*$, alors (D) admet une solution optimale $\lambda^*$ et leurs valeurs sont égales : $c^T x^* = b^T \lambda^*$.
\end{itemize}

\section{Géométrie de l'ensemble réalisable}

L'ensemble des solutions réalisables $\mathcal{F} = \{ x \in \mathbb{R}^n \mid Ax = b, x \geq 0 \}$ est un polyèdre convexe (intersection d'un sous-espace affine et de l'orthant positif).

\begin{proposition}
    Les sommets (ou points extrêmes) de ce polyèdre correspondent exactement aux solutions de base réalisables du problème linéaire.
\end{proposition}

La méthode du simplexe consiste géométriquement à se déplacer d'un sommet à un sommet adjacent du polyèdre, en suivant une arête qui améliore (diminue) la fonction objectif.
Puisque le nombre de sommets est fini (au plus $\binom{n}{m}$), l'algorithme termine en un nombre fini d'étapes (sous réserve de précautions contre le cyclage, par exemple avec la règle de Bland).

\section{Aperçu de la méthode du simplexe}

\begin{algorithm}[H]
    \caption{One step of the Simplex method}
    \SetAlgoLined
    Given $\mathcal{B}$, $\mathcal{N}$, $x_B=B^{-1}b\geq 0$, $x_N=0$\;
    Solve $B^\top\lambda = c_B$ for $\lambda$\;
    Compute $s_N=c_N - N^\top\lambda$\; \Comment{Pricing}
    \If{$s_N\geq 0$}{
        \Stop \Comment{optimal point found}
    }
    Select $q\in\mathcal{N}$ with $s_q<0$ as the entering index\;
    Solve $Bd=A_q$ for $d$\;
    \If{$d\leq 0$}{
        \Stop \Comment{problem is unbounded}
    }
    Calculate $x_q^+=\min_{i \mid d_i> 0}\left(\frac{(x_B)_i}{d_i}\right)$, use $p$ to denote the minimizing $i$\;
    Update $x_B^+=x_B - d x_q^+$, $x_N^+=(0, \ldots, 0, x_q^+, 0, \ldots, 0)^\top$\;
    Change $\mathcal{B}$ by adding $q$ and removing the basic variable corresponding to column $p$ of $B$\;
\end{algorithm}
